\documentclass[12pt]{article}

% =========================================================
% PAQUETES
% =========================================================
\usepackage[spanish]{babel}
\usepackage[utf8]{inputenc}
\usepackage[T1]{fontenc}
\usepackage{lmodern}
\usepackage{amsmath}
\usepackage{geometry}
\usepackage{array}
\usepackage{booktabs}
\usepackage{multirow}
\usepackage{xcolor}

\geometry{letterpaper, margin=2cm}

\setlength{\parskip}{0.4em}
\setlength{\parindent}{0pt}

% =========================================================
% DOCUMENTO
% =========================================================
\begin{document}
	
	\begin{center}
		{\Large \textbf{PAUTA DE EVALUACIÓN UNIDAD 1 (VIDEO)}}\\[0.5cm]
	\end{center}
	
	\textbf{NOMBRES DE INTEGRANTES:}\\[0.3cm]
	\hrulefill\\[0.3cm]
	\hrulefill\\[0.3cm]
	\hrulefill
	
	\vspace{0.5cm}
	
	% =========================================================
	\section*{Instrucciones Generales}
	
	\begin{itemize}
		\item El video trata los temas de la Unidad I.
		\item El video tiene una duración máxima de 16 minutos.
		\item El video debe considerar la explicación y desarrollo completo de cada problema.
		\item Todos los integrantes deben presentarse en forma visible en el video.
	\end{itemize}
	
	\vspace{0.3cm}
	
	\textbf{Puntaje:}
	
	\begin{center}
		\begin{tabular}{c c c}
			\textbf{Conocimiento 90\%} & \textbf{Aspectos Generales 5\%} & \textbf{Proceso 5\%}
		\end{tabular}
	\end{center}
	
	\vspace{0.3cm}
	
	\textbf{Escala de evaluación:}
	
	\begin{center}
		\begin{tabular}{|c|c|}
			\hline
			0 & No se cumple \\
			\hline
			1 & Se cumple en lo mínimo aceptado \\
			\hline
			2 & Se cumple de manera aceptable \\
			\hline
			3 & Se cumple de manera sobresaliente \\
			\hline
		\end{tabular}
	\end{center}
	
	% =========================================================
	\section*{Rúbrica de Evaluación}
	
	\renewcommand{\arraystretch}{1.5}
	
	\begin{tabular}{|m{3cm}|m{10cm}|m{2cm}|}
		\hline
		\textbf{Categoría} & \textbf{Aspectos a evaluar} & \textbf{Puntaje} \\
		\hline
		
		\multirow{5}{*}{Aspectos Generales (5\%)} 
		& El video tiene portada con el nombre de todos los integrantes del grupo. & \\
		\cline{2-3}
		& La calidad del audio e imagen es excelente. & \\
		\cline{2-3}
		& El video cumple con el tiempo definido (16 minutos máximo). & \\
		\cline{2-3}
		& Se cumple el plazo de entrega del video. & \\
		\cline{2-3}
		& Cada estudiante aparece en la presentación del problema. & \\
		\hline
		
		\multirow{3}{*}{Proceso (5\%)} 
		& Evaluación de proceso en ayudantía. & \\
		\cline{2-3}
		& Autoevaluación entregada en Aula Virtual. & \\
		\cline{2-3}
		& Coevaluación entregada en Aula Virtual. & \\
		\hline
		
		\multirow{4}{*}{Aplicación (40\%)} 
		& Desarrolla correctamente la resolución del problema 1 en su totalidad. & \\
		\cline{2-3}
		& Desarrolla correctamente la resolución del problema 2 en su totalidad. & \\
		\cline{2-3}
		& Desarrolla correctamente la resolución del problema 3 en su totalidad. & \\
		\cline{2-3}
		& Desarrolla correctamente la resolución del problema 4 en su totalidad. & \\
		\hline
		
		\multirow{4}{*}{Dominio (50\%)} 
		& Explica correctamente la resolución del problema 1. Demuestra dominio del tema sin leer ni apoyarse en apuntes. & \\
		\cline{2-3}
		& Explica correctamente la resolución del problema 2. Demuestra dominio del tema sin leer ni apoyarse en apuntes. & \\
		\cline{2-3}
		& Explica correctamente la resolución del problema 3. Demuestra dominio del tema sin leer ni apoyarse en apuntes. & \\
		\cline{2-3}
		& Explica correctamente la resolución del problema 4. Demuestra dominio del tema sin leer ni apoyarse en apuntes. & \\
		\hline
		
	\end{tabular}
	
	\vspace{0.8cm}
	
	\textbf{Observaciones:}
	
	\vspace{2cm}
	
	\hrulefill
	
\end{document}